% 解説プリントのコメント（saiten が書く。./tools/kaisetsu.ps1 がこのファイルからプリントを作る）
% 量子力学1-6（測定と不確定性関係）2026-11-01 のスキャン（問題1〜6）
%% matome
\begin{daiji}
  \item 交換子の計算で，負号をかけた括弧を外すときは括弧の中のすべての項の符号を変える。最後に関数 $\psi$ に作用させて直接計算し，係数と符号を検算する（問題2 (4)）。
  \item 「交換子が $0$ でないから同時に定められない」で止めない。同時固有関数 $\phi$ を仮定すると交換子を作用させて $0$ になり，交換子が $0$ でない定数なら $\phi=0$ となって矛盾する，という筋を書く（問題3 (3)。問題5 (4) も同じ型）。
  \item 問われたものをそのまま答える。「得られる値と確率」「エネルギーの確率の比較」を期待値で代用しない（問題1 (1)(4)）。数値の答えには正しい単位を付ける（問題6 (5)）。
\end{daiji}
\begin{maemo}
  \item 交換子の符号：量子力学1-2 問題4 (2)（2026-10-14）で交換子の符号を逆（$-3\ii\hbar$）にした。今回も問題2 (4) で符号を誤った。交換子を出したら，必ず $\psi$ に作用させる検算を1行入れる。
  \item 理由の書き落とし：量子力学1-5 問題7 (4)（2026-10-30）で，成り立たない段階を示しただけで理由を書かなかった。今回も問題3 (3) で「$\ne 0$ より」で止まった。「何がどう矛盾するか」を1行で書く。
\end{maemo}
%% mondai 1
\begin{kotaeran}
  \machigai{1}{○}{記述} 規格化の確認 $\int\abs{\psi}^{2}\dd x=(1/25)\int(9\abs{\phi_{1}}^{2}+16\abs{\phi_{2}}^{2})\dd x=1$ は正しい。問われている「得られる値とその確率」を答えとして書いていない（(2) の $\expval{E}=(9/25)E_{1}+(16/25)E_{2}$ の重みには出ている）。答えは $E_{1}$（確率 $9/25$）と $E_{2}=4E_{1}$（確率 $16/25$）。
  \machigai{4}{○}{記述} $\expval{x}$ は $\psi$ で $L/2$，$\chi$ で $\approx\num{0.327}L$ と正しい。エネルギーについては「$\expval{E}_{\chi}=(73/25)E_{1}$ で同じ」と期待値を比べたが，問われているのは確率である。期待値が同じでも確率が同じとは限らない。$\chi$ の係数 $3/5$，$4/5$ の絶対値は $\psi$ の $3/5$，$4\ii/5$ と同じなので，確率（$E_{1}$ が $9/25$，$4E_{1}$ が $16/25$）が同じ，と書く。なお途中の式で $L/2-128L/(75\pi^{2})$ を $128L/(75\pi^{2})$ と書いたが，数値 $\num{0.327}L$ は正しい（減点なし）。
\end{kotaeran}
\begin{rikai}
  \item 測定で得られる値は固有値だけで，その確率は展開係数の絶対値の2乗 $\abs{c_{n}}^{2}$ である。期待値 $\sum_{n}a_{n}\abs{c_{n}}^{2}$ は確率から決まるが，逆に期待値だけからは確率は決まらない。
  \item 係数の位相（$4\ii$ か $4$ か）は確率を変えないが，交差項 $2\Re(c_{1}^{*}c_{2})X_{12}$ を通して $\expval{x}$ を変える。$\psi$ では $c_{1}^{*}c_{2}$ が純虚数なので交差項が消え，$\chi$ では残る。
\end{rikai}
%% mondai 2
\begin{kotaeran}
  \machigai{4}{△}{計算} $\comm{\hat{p}}{\hat{x}^{3}}=-\comm{\hat{x}\hat{x}^{2}}{\hat{p}}=-\left(\hat{x}\comm{\hat{x}^{2}}{\hat{p}}+\comm{\hat{x}}{\hat{p}}\hat{x}^{2}\right)$ までは正しい。次の行を $-2\ii\hbar\hat{x}^{2}+\ii\hbar\hat{x}^{2}$ と書き，括弧の2項目の符号を変えていない。正しくは $-2\ii\hbar\hat{x}^{2}-\ii\hbar\hat{x}^{2}=-3\ii\hbar\hat{x}^{2}$。
\end{kotaeran}
\begin{rikai}
  \item 負号を前に出した括弧を外すときは，括弧の中のすべての項の符号を変える。
  \item 検算：$\psi$ に作用させて $\hat{p}(x^{3}\psi)-x^{3}\hat{p}\psi=-\ii\hbar(3x^{2}\psi+x^{3}\psi')+\ii\hbar x^{3}\psi'=-3\ii\hbar x^{2}\psi$ を直接計算する。問題3 (1) の $\comm{V}{\hat{p}}=\ii\hbar V'(x)$ で $V=x^{3}$ とおいても $\comm{\hat{p}}{\hat{x}^{3}}=-3\ii\hbar\hat{x}^{2}$ が出る（$x^{n}$ の微分と同じ係数 $n$ が付く）。
\end{rikai}
%% mondai 3
\begin{kotaeran}
  \machigai{3}{○}{記述} $\comm{\hat{H}}{\hat{p}}=-\ii\hbar F\ne 0$ は正しい。理由が「$\ne 0$ より同時観測不可能」で止まっている。交換子が $0$ でないだけでは，同時固有関数が1つも無いとは言えない（問題5 の $A$，$B$ は交換しないが，$\vect{e}_{1}=(1,0,0)^{\mathrm{T}}$ は両方の固有ベクトル）。同時固有関数 $\phi$（$\hat{H}\phi=E\phi$，$\hat{p}\phi=p\phi$）があったとすると $\comm{\hat{H}}{\hat{p}}\phi=pE\phi-Ep\phi=0$，一方で左辺は $-\ii\hbar F\phi$ で $F\ne 0$ なので $\phi=0$ となり，固有関数であることに反する，まで書く。
\end{kotaeran}
\begin{rikai}
  \item 要点は交換子が $0$ でない「定数」であること：$0$ でない定数をかけて $0$ になる関数は $0$ しか無いので，同時固有関数は1つも無い。交換子が演算子（行列）のときは，作用させて $0$ になる特定の状態がありうる。
\end{rikai}
%% mondai 5
\begin{kotaeran}
  \machigai{4}{○}{計算} $AB$，$BA$ は正しい。$AB-BA$ の $(3,2)$ 成分を $2$ と書いたが，$(-1)-1=-2$ である（$\ne 0$ という結論は変わらない）。理由は「$B$ の測定で $A=-1$ となる状態が含まれる」まで書けている。その前に，$AB\ne BA$ なので $A$，$B$ の両方の固有ベクトルだけで規格直交な3本の組はつくれない（つくれれば $AB=BA$ になる），だから $B$ の固有ベクトル $\vect{w}$ が $A$ の固有ベクトルとは限らない，という論理をつなぐ。
\end{kotaeran}
\begin{rikai}
  \item 検算：$A$，$B$ が実対称行列なら $(AB-BA)^{\mathrm{T}}=BA-AB$ なので，$AB-BA$ は反対称行列になる。対角成分が $0$ で，$(2,3)$ 成分と $(3,2)$ 成分が符号違いかを確かめる。
  \item 同時固有ベクトルの規格直交な組があれば，その各ベクトルで $AB\vect{y}=BA\vect{y}$ となり，任意のベクトルを展開して $AB=BA$ になる。対偶から，$AB\ne BA$ なら $B$ の測定後の状態が $A$ の固有ベクトルになるとは限らず，最初の $A$ の結果が保たれない。
\end{rikai}
%% mondai 6
\begin{kotaeran}
  \machigai{5}{○}{単位} $\delta p$ の数値 $\num{7.46e-25}$ は正しいが，単位を \unit{\metre} と書いた。運動量の単位は \unit{\kilogram\metre\per\second} である。$\delta v$ と平均の速さの単位 \unit{\metre\per\second} は正しい。
\end{kotaeran}
\begin{rikai}
  \item 数値に単位を付けるときは，式の次元から出す：$\hbar/\sigma$ は \unit{\joule\second}/\unit{\metre}$=$\unit{\kilogram\metre\per\second} で運動量の単位になる。
\end{rikai}
